\section{Introduction}
\label{sec:introduction}
The convergence of research and design endeavours in architecture and interaction design (IxD) has brought forward new ways of thinking about the relationship between humans and the built environment. Historically, research in the domains of architecture and building performance engineering has focused on human experiences through an objective lens \cite{kirsh2019architects}, often using environmental sensing and metrics such as thermal comfort and indoor air quality to assess and improve building design \cite{chappells2004comfort}.
In human-computer interaction (HCI), questions around architectural spaces began to emerge, notably in the foundational paper on \emph{places} and \emph{spaces} by \citet{harrison1996re} in \citeyear{harrison1996re}, which explored how people interact with physical spaces through a digital lens \cite{harrison1996re}. 
Across various domains, researchers have since developed a deeper understanding of how technology shapes human experiences within built environments, moving beyond objective assessments of the environment to recognise individual differences and subjective experiences. In recent years, researchers and designers have sought to re-imagine the role of interactive technologies embedded in our built environments  \cite{wiberg2020interaction, alavi2018artifacts}. This has been complemented by new visions that interface HCI with the domains of architecture and urban design to delineate a new area of design research under the banner of human-building interaction (HBI) \cite{alavi2019introduction}. HBI delves into the relationships between digital technologies, physical spaces, and the human experiences they shape \cite{alavi2019introduction,becerik2022field}, considering the subjectivity and individuality of human experiences. 

Furthermore, HBI acknowledges that technologies impact human experiences in smart building environments in various ways. Designers and developers who create and shape smart buildings inherently embed both explicit and implicit values, assumptions, and expectations from their respective disciplines into the experiences they envision or the technologies they create~\cite{wiberg2020interaction}. For example, they might focus primarily on enhancing human comfort (across its four classic dimensions: thermal, visual, acoustic, respiratory) \cite{alavi2017comfort}, or on human agency and control \cite{brambilla2017our}.
They might promote supporting individual needs \cite{baillie2008place}, or collaborative, community-centred interactions~\cite{jacobsen2023living}. They might frame the role of the built environment in people's lives purely as functional spaces \cite{yang2014making}, or as immersive, responsive environments that contribute to emotions \cite{mcdonald2015nature}, sustainability \cite{loh2010please}, and social engagement \cite{schuricht2007freequent}. Furthermore, with the rapid and pervasive deployment of interactive technology, smart building environments are increasing in complexity, and the factors involved in shaping human experiences are also increasingly multifaceted. By and large, smart building research consists of a rich but fragmented body of work encompassing diverse perspectives, methodologies, and disciplines. This highlights the necessity to conduct a comprehensive mapping of the literature in order to understand the full scope of what we design for when we design smart buildings. Our review studies human experiences in smart buildings through this central question: \textit{what do we envision and research when we design for human experiences in future smart buildings?} 

The answers to this question not only shape how smart building experiences are created, but also how people interact with, value, and experience the spaces they inhabit. Smart buildings, in this perspective, are not just passive infrastructures; they are active ``shapers'' of individual, cultural, and societal attitudes toward the built environment. This paper seeks to uncover these visions through a scoping review, addressing the following research questions:
\begin{itemize}
    \item \textbf{\textit{RQ1.}} What are the different types of \emph{human experiences} that research in the broad domain of smart buildings aims to improve or introduce?
    \item \textbf{\textit{RQ2.}} What are the common \emph{design methods and approaches} employed in smart building research to shape human experiences? 
    \item \textbf{\textit{RQ3.}} What types of \emph{technological interventions} are developed or conceptualised to enhance interactive experiences in buildings?
\end{itemize}



To address these questions, we analysed 192 papers that contribute to the research and design of human experiences in smart buildings, sourced from the ACM Digital Library and Scopus. Using a hybrid latent coding approach \cite{braun2023thematic}, our scoping review provides a comprehensive descriptive analysis of how smart building research has explored human experiences in relation to digital technology. Specifically, we offer: (a)\emph{a classification of 11 distinct human experiences} that smart building research aims to improve or introduce, (b) \emph{a classification of 20 commonly used design mechanisms} employed to enhance or create these human experiences, and (c) \emph{a typology of the technological interventions} used to impact these human experiences.


Our contributions provide a foundation for understanding and describing the current and future of smart building design research, while incorporating a broad, diverse range of human experience in these spaces.  
